\section{Unprespecified and Exploratory Analyses}

\subsection{Analyses Not Specified in the Primary Analysis Plan}

Analyses that are not specified in the preregistered analysis plan may be conducted where they provide useful additional information about characteristics, behaviours, or patterns observed within the survey dataset. Such analyses will be clearly distinguished from prespecified analyses and will not be presented as having been planned in advance.

\subsection{Investigation of Unexpected or Unanticipated Patterns}

\sourceimage{1510748}{Nopony expects the Ponkquisition!}

Unexpected or unanticipated patterns identified during analysis may be investigated where they are substantively interesting or relevant to the study population. Such investigations will be treated as exploratory and will not be used to retrospectively redefine the primary research questions.

\subsection{Data-driven and Additional Analyses}

Additional data-driven analyses may include cross-tabulations, subgroup analyses, association analyses, clustering, or other descriptive approaches appropriate to the structure of the dataset. Analyses arising from observed findings may examine additional variables, alternative subgroup definitions, or more detailed intersections.

Any analysis motivated by patterns observed in the data will be identified as exploratory and interpreted with appropriate caution.

% TBD: The specific types of unprespecified analysis permitted, including whether more advanced statistical or
% computational approaches will be used, will be established before analysis where practicable.

\subsection{Classification and Reporting of Exploratory Analyses}

An analysis will be classified as \textbf{exploratory} if it was not sufficiently specified in the preregistered analysis plan before examination of the substantive study data, or if its inclusion, specification, subgroup definition, or analytical approach was materially motivated by patterns observed in the data.

Exploratory analyses will be clearly distinguished from prespecified analyses in the final publication. They may be used to generate potential research questions or hypotheses for subsequent studies, but will not be presented as confirmatory evidence for hypotheses formulated after observing the data.

\subsection{Deviations from the Preregistered Plan}

All substantive deviations from the preregistered analysis plan will be documented when identified and retained alongside the original preregistration and analysis materials. The original preregistered plan will not be retrospectively altered in a manner that obscures decisions specified before data analysis.

Analyses will be classified according to when and why they were specified. \textbf{Preregistered analyses} are those specified before examination of the substantive study data. \textbf{Modified analyses} are analyses specified in the preregistration but subsequently changed in a substantive way. \textbf{Subsequently added analyses} are analyses not included in the preregistered plan and introduced after registration, including analyses motivated by observations made during analysis.

Material deviations will be reported in the final research paper and, where appropriate, in supplementary materials or the associated research repository. The report will identify what was originally planned, what was subsequently changed, and whether the change affected the interpretation or scope of the resulting analysis.

Where a deviation occurs, its rationale will be documented. Deviations will not be justified solely on the basis that a revised analysis produces a more desirable or statistically significant result. Where a change is motivated by an unexpected finding in the data, this will be explicitly identified and the resulting analysis will generally be classified as exploratory.

Minor corrections or implementation details that do not materially alter the planned analysis may be documented separately from substantive deviations.

% TBD: The documentation system, threshold for reportable substantive deviations, location and format of deviation
% reporting, and precise rationale requirements will be established before analysis.
